\introChapter{ความสัมพันธ์และฟังก์ชัน (Relations and Functions)}

\begin{description}[leftmargin=2cm, style=nextline, itemsep=0.1em, parsep=0pt]
    \item[รหัสวิชา] 4091606
    \item[ชื่อวิชา] คณิตศาสตร์สำหรับคอมพิวเตอร์
    \item[หน่วยกิต] 3 (3-0-6)
    \item[บทที่] 3
    \item[ชื่อบท] ความสัมพันธ์และฟังก์ชัน (Relations and Functions)
    \item[จำนวนชั่วโมง] 9 ชั่วโมง
\end{description}

\noindent \textbf{\large วัตถุประสงค์การเรียนรู้ประจำบท}\\[0.2em]
\noindent เมื่อนักศึกษาเรียนบทนี้จบแล้ว นักศึกษาจะสามารถ
\begin{enumerate}[topsep=0.2em, itemsep=0.15em, leftmargin=*]
    \item อธิบายความหมายของความสัมพันธ์ และผลคูณคาร์ทีเซียนได้
    \item ระบุสมบัติของความสัมพันธ์ เช่น สะท้อน สมมาตร ปฏิสมมาตร ถ่ายทอด และการปิดสมบัติได้
    \item ระบุและจำแนกประเภทของฟังก์ชัน เช่น ฟังก์ชันหนึ่งต่อหนึ่ง ฟังก์ชันทั่วถึง ฟังก์ชันหนึ่งต่อหนึ่งและทั่วถึง ได้
    \item คำนวณลำดับเลขคณิต ลำดับเรขาคณิต อนุกรมเลขคณิต และอนุกรมเรขาคณิตได้
    \item ประยุกต์ใช้ความสัมพันธ์และฟังก์ชันในบริบทของโปรแกรมคอมพิวเตอร์ได้
\end{enumerate}

\noindent \textbf{\large หัวข้อที่สอน}\\[0.2em]
\begin{enumerate}[topsep=0.2em, itemsep=0.15em, leftmargin=*]
    \item พื้นฐานความสัมพันธ์และคู่อันดับ
    \item สมบัติและการขยายความสัมพันธ์ (สมบัติ 4 ประการ และการปิดสมบัติ)
    \item ฟังก์ชัน (นิยาม โดเมน/เรนจ์ ประเภท การประกอบ ฟังก์ชันผกผัน และฟังก์ชันพิเศษ)
    \item ลำดับและอนุกรม (ลำดับเลขคณิต, ลำดับเรขาคณิต, อนุกรมเลขคณิต, อนุกรมเรขาคณิต)
    \item การประยุกต์ใช้ในวิทยาการคอมพิวเตอร์
\end{enumerate}

\noindent \textbf{\large สื่อการสอน}
\begin{enumerate}[topsep=0.2em, itemsep=0.15em, leftmargin=*]
    \item สไลด์ประกอบการสอน
    \item ตัวอย่างการคำนวณและแผนภาพ
    \item แบบฝึกหัดท้ายบท
\end{enumerate}
